\documentclass{article}

\usepackage{fullpage}
\usepackage{url}
\usepackage{graphicx}
\usepackage{amsmath}
\usepackage{physics}
\usepackage{mathtools}
\usepackage{bbold}

\DeclareMathOperator{\erfi}{erfi}
\newcommand{\R}{\ensuremath{\mathbb{R}}}
\newcommand{\C}{\ensuremath{\mathbb{C}}}

\title{Differential Geometry Excercise 4}
\author{Idan Stark}
\date{\today}

\begin{document}
\maketitle

\section*{Question 1}
Let $M \subset \R^3$ be a two-dimensional, embedded submanifold with the induced Riemannian metric.
Let $U \subset \R^3$ be a neighborhood of $0$ and let $f: \R^2 \rightarrow \R$ be a smooth function with $f(0) = 0$ and $\nabla f(0) = 0$. Let us assume that
\begin{equation}
    M \cap U = \left\{(x,y,z) \in U : z = f(x,y)\right\}
\end{equation}
\subsection*{Local orthonormal coframe}
The induced metric comes from $\R^3$ Eucledian metric:
\begin{equation}
    g_{\R^3} = dx^2 + dy^2 + dz^2
\end{equation}
Inside $M \cap U$, the induced metric takes the form:
\begin{equation}
    g_{M \cap U} = dx^2 + dy^2 + \left(f_x dx + f_y dy\right)^2 =
    \left(1 + f_x^2\right)dx^2 + 2f_xf_ydxdy + \left(1+ f_y^2\right)dy^2
\end{equation}
Where
\begin{equation}
    f_x = \frac{\partial f}{\partial x}
    \qquad
    f_y = \frac{\partial f}{\partial y}
\end{equation}
To find a coframe, let us use Cholesky decomposition. Define the two one-forms:
\begin{equation}
    \omega_1 = a_x dx
    \qquad
    \omega_2 = b_x dx + b_y dy
\end{equation}
Such that:
\begin{equation}
    g = \omega_1^2 + \omega_2^2
\end{equation}
Openning this up gives:
\begin{equation}
    g = (a_x^2 + b_x^2) dx^2 + 2b_xb_ydxdy + b_y^2dy^2
\end{equation}
comparing this to the original metric gives:
\begin{equation}
\begin{gathered}
    a_x = \sqrt{1 + \frac{f_x^2}{1 + f_y^2}}
    \qquad
    b_x = \frac{f_xf_y}{\sqrt{1 + f_y^2}}
    \qquad
    b_y = \sqrt{1 + f_y^2}
\end{gathered}
\end{equation}
And so, and orthonotmal frame is:
\begin{equation}
    \omega_1 = \sqrt{1 + \frac{f_x^2}{1 + f_y^2}}dx
    \qquad
    \omega_2 = \frac{f_xf_y}{\sqrt{1 + f_y^2}}dx + \sqrt{1 + f_y^2}dy
\end{equation}
\subsection*{One form derivation}
Let
\begin{equation}
    \omega = \omega_1 + i\omega_2
\end{equation}
And let us look for the 1-form $\varphi$ such that $d\omega = i\varphi \wedge \omega$. Plugging the defintion of $\omega$ into it gives:
\begin{equation}
    d\omega_1 + id\omega_2 = i \varphi \wedge (\omega_1 + i\omega_2) = -\varphi \wedge \omega_2 + i \varphi \wedge \omega_1
\end{equation}
Which requires:
\begin{equation}
\begin{cases}
    d\omega_1 = -\varphi \wedge \omega_2
    \\
    d\omega_2 = \varphi \wedge \omega_1
\end{cases}
\end{equation}
Since $\omega_1, \omega_2$ is a coframe, $\varphi$ can be written as its linear combination $\varphi = \varphi_1 \omega_1 + \varphi_2\omega_2$:
\begin{equation}
\begin{cases}
    d\omega_1 = -\varphi_1 \omega_1 \wedge \omega_2
    \\
    d\omega_2 = - \varphi_2 \omega_1 \wedge \omega_2
\end{cases}
\end{equation}
Let us now derive the values of these expressions in terms of $dx$ and $dy$:
\begin{equation}
    \omega_1 \wedge \omega_2 = \sqrt{\left(1 + \frac{f_x^2}{1 + f_y^2}\right)\left(1 + f_y^2\right)} dx \wedge dy = 
    \sqrt{1 + f_x^2 + f_y^2} dx \wedge dy
\end{equation}
\begin{equation}
\begin{gathered}
    d\omega_1 = -\frac{\partial a_x}{\partial y} dx \wedge dy =
    -a_x^{-1}\left(\frac{(1 + f_y^2)f_xf_{xy} - f_x^2f_yf_{yy}}{\left(1 + f_y^2\right)^2}\right) dx \wedge dy = \\ =
    -\frac{f_xf_{xy} + f_y^2f_xf_{xy} - f_x^2f_yf_{yy}}{\left(1 + f_y^2\right)^{3/2}\left(1 + f_x^2 + f_y^2\right)^{1/2}} dx \wedge dy = \\ =
    -\frac{f_x(f_{xy} + f_y^2f_{xy} - f_xf_yf_{yy})}{\left(1 + f_y^2\right)^{3/2}\left(1 + f_x^2 + f_y^2\right)^{1/2}} dx \wedge dy
\end{gathered}
\end{equation}
\begin{equation}
\begin{gathered}
    d\omega_2 = \left(\frac{\partial b_y}{\partial x} - \frac{\partial b_x}{\partial y}\right) dx \wedge dy = \\ =
    \left(
        \frac{f_yf_{xy}}{\sqrt{1 + f_y^2}} -
        \frac{(f_{xy}f_y + f_xf_{yy})(1 + f_y^2) - f_xf_y^2f_{xy}}{(1 + f_y^2)^{3/2}}
    \right) dx \wedge dy = \\ =
    \left(
        \frac{(f_{xy}f_y - f_xf_{yy} - f_yf_{xy})(1 + f_y^2) + f_xf_y^2f_{xy}}{(1 + f_y^2)^{3/2}}
    \right) dx \wedge dy = \\ =
    -\frac{f_x(f_{yy} + f_y^2f_{yy} + f_y^2f_{xy})}{(1 + f_y^2)^{3/2}}
    dx \wedge dy
\end{gathered}
\end{equation}
This gives:
\begin{equation}
    \varphi_1 = \frac{f_x(f_{xy} + f_y^2f_{xy} - f_xf_yf_{yy})}{\left(1 + f_y^2\right)^{3/2}\left(1 + f_x^2 + f_y^2\right)}
    \qquad
    \varphi_2 = \frac{f_x(f_{yy} + f_y^2f_{yy} + f_y^2f_{xy})}{(1 + f_y^2)^{3/2}\left(1 + f_x^2 + f_y^2\right)^{1/2}}
\end{equation}

\subsection*{Theorema Egregium}
Let us now proove Theorema Egregium at the origin, i.e. show that:
\begin{equation}
    d\varphi = K \omega_1 \wedge \omega_2 \qquad s.t. \qquad K = \det \nabla^2 f(0)
\end{equation}
To show that, let us expand $f_x$ and $f_y$ around the origin (all second derivatives from this point on are calculated at the origin, i.e. $f_{xx} = f_{xx}(0)$ and so forth):
\begin{equation}
    f_x \approx f_{xx}x + f_{xy}y
    \qquad
    f_y \approx f_{xy}x + f_{yy}y
\end{equation}
Pluggin that into the formulas for $\varphi_1$ and $\varphi_2$ and keeping only linear terms:
\begin{equation}
    \varphi_1 = \frac{f_x(f_{xy} + f_y^2f_{xy} - f_xf_yf_{yy})}{\left(1 + f_y^2\right)^{3/2}\left(1 + f_x^2 + f_y^2\right)} \approx
    f_{xy}(f_{xx}x + f_{xy}y)
\end{equation} 
\begin{equation}
    \varphi_2 = \frac{f_x(f_{yy} + f_y^2f_{yy} + f_y^2f_{xy})}{(1 + f_y^2)^{3/2}\left(1 + f_x^2 + f_y^2\right)^{1/2}} \approx
    f_{yy}(f_{xx}x + f_{xy}y)
\end{equation}
Which gives the form:
\begin{equation}
\begin{gathered}
    \varphi = f_{xy}(f_{xx}x + f_{xy}y) \omega_1 + f_{yy}(f_{xx}x + f_{xy}y) \omega_2 = \\ =
    (f_{xx}x + f_{xy}y)\left(f_{yy} \frac{f_xf_y}{\sqrt{1 + f_y^2}} + f_{xy} \sqrt{1 + \frac{f_x^2}{1 + f_y^2}}\right)dx + f_{yy}(f_{xx}x + f_{xy}y)\sqrt{1 + f_y^2}dy \approx \\ \approx
    f_{xy}(f_{xx}x + f_{xy}y)dx + f_{yy}(f_{xx}x + f_{xy}y)dy
\end{gathered}
\end{equation}
Deriving this, we get:
\begin{equation}
\begin{gathered}
    d\varphi = \left(
        \frac{\partial}{\partial x}\left(f_{yy}(f_{xx}x + f_{xy}y)\right) -
        \frac{\partial}{\partial y}\left(f_{xy}(f_{xx}x + f_{xy}y)\right)
    \right) dx \wedge dy = \\ =
    \left(f_{xx}f_{yy} - f_{xy}^2\right) dx \wedge dy = \det \nabla^2 f(0) dx \wedge dy
\end{gathered}
\end{equation}
But, in the origin:
\begin{equation}
    \omega_1 \wedge \omega_2 = \sqrt{1 + f_x^2 + f_y^2} dx \wedge dy = dx \wedge dy
\end{equation}
Thus:
\begin{equation}
    d\varphi = \det \nabla^2 f(0) \omega_1 \wedge \omega_2
\end{equation}

\section*{Question 2}
\subsection*{The gaussian curvature in $\mathbb{H}^2$}
Let $\mathbb{H}^2$ be the upper half plane with the metric:
\begin{equation}
    g = \frac{dx^2 + dy^2}{y^2}
\end{equation}
Let us do Gram-Schmidt to find an orthonormal frame $E_1, E_2$. Starting from
\begin{equation}
    E_1 = \partial_x
\end{equation}
And then:
\begin{equation}
    E_2 = \partial_y - \frac{y^{-2}\left(0\cdot1 + 1\cdot0\right)}{y^{-2}}\partial_x = \partial_y
\end{equation}
So $\partial_x, \partial_y$ are already an orthogonal frame. Normalizing it gives:
\begin{equation}
    E_1 = y\partial_x \qquad E_2 = y\partial_y    
\end{equation}
and so and orthonormal coframe is:
\begin{equation}
    \omega_1 = y^{-1}dx \qquad \omega_2 = y^{-1}dy
\end{equation}
And so
\begin{equation}
    d\omega_1 = y^{-2}dx \wedge dy
    \qquad
    d\omega_2 = 0
\end{equation}
So we can concluate that:
\begin{equation}
    \varphi = -\omega_1
\end{equation}
Since then:
\begin{equation}
    d\omega_1 = -\varphi \wedge \omega_2 = \omega_1 \wedge \omega_2 = y^{-2}dx \wedge dy
    \qquad
    d\omega_2 = \varphi \wedge \omega_1 = -\omega_1 \wedge \omega_1 = 0
\end{equation}
Now, since
\begin{equation}
    \omega_1 \wedge \omega_2 = y^{-2} dx \wedge dy
\end{equation}
then:
\begin{equation}
    d\varphi = -d\omega_1 = -y^{-2}dx \wedge dy = -\omega_1 \wedge \omega_2
\end{equation}
And so:
\begin{equation}
    K = -1
\end{equation}
\subsection*{The Sphere $S^2$}
For the sphere. Using the stereographic coordinate system, the metric is:
\begin{equation}
    g = \frac{4}{(1 + x^2 + y^2)^2}\left(dx^2 + dy^2\right)
\end{equation}
Similar to the half-plane, this gives the orthonormal frame:
\begin{equation}
    E_1 = \frac{1 + x^2 + y^2}{2}\partial_x
    \qquad
    E_2 = \frac{1 + x^2 + y^2}{2}\partial_y
\end{equation}
And the orthonormal coframe:
\begin{equation}
    \omega_1 = \frac{2}{1 + x^2 + y^2}dx
    \qquad
    \omega_2 = \frac{2}{1 + x^2 + y^2}dy
\end{equation}
the derivatives are then:
\begin{equation}
\begin{gathered}
    d\omega_1 = -\frac{4ydy}{\left(1 + x^2 + y^2\right)^2}\wedge dx = \frac{4ydx}{\left(1 + x^2 + y^2\right)^2}\wedge dy = \frac{2y}{1 + x^2 + y^2} dx \wedge d\omega_2
    \\
    d\omega_2 = -\frac{4xdx}{\left(1 + x^2 + y^2\right)^2} \wedge dy = \frac{4xdy}{\left(1 + x^2 + y^2\right)^2}\wedge dx = \frac{2x}{1 + x^2 + y^2} dy \wedge d\omega_1
\end{gathered}
\end{equation}
And the cannoincal connection form is:
\begin{equation}
    \varphi = \frac{2xdy - 2ydx}{1 + x^2 + y^2}
\end{equation}
Deriving it:
\begin{equation}
    d\varphi = \frac{4}{1 + x^2 + y^2} dx \wedge dy
\end{equation}
On the other hand
\begin{equation}
    \omega_1 \wedge \omega_2 = \frac{4}{1 + x^2 + y^2} dx \wedge dy
\end{equation}
So
\begin{equation}
    K = 1
\end{equation}
\subsection*{The Flat Torus}
Let us look at the torus with the flat metric $T = \R^2 / \mathbb{Z}^2 $. This Torus has the same metric as $\R^2$, and is therefore flat. Let us show it. Starting from the metric:
\begin{equation}
    g = dx^2 + dy^2
\end{equation}
This gives:
\begin{equation}
    \omega_1 = dx \qquad \omega_2 = dy
\end{equation}
Then:
\begin{equation}
    d\omega_1 = d\omega_2 = 0
\end{equation}
And
\begin{equation}
    \varphi = 0
\end{equation}
Therefore:
\begin{equation}
    K = 0
\end{equation}
\section*{Question 3}
Let $M$ be a smooth manifold. Let $\left\{U_i\right\}$ be a cover of $M$ and $\left\{\psi\right\}$ a partition of unity over $M$ subordinate to this cover. Then, for every subset $U_i$ define a smooth orthonormal coframe $\omega_1,\dots,\omega_n$ on $U_i$. This gives a locally defined (0,2)-tensor field:
\begin{equation}
    g_i = \sum_j^n \omega_j \otimes \omega_j
\end{equation}
that is symmetric and positive. The global (0,2)-tensor field
\begin{equation}
    g = \sum_i \psi_i g_i
\end{equation}
is smooth and thus a metric.
\newline
Now, let $(p, X) \in TM$ be a member of the tangent bundle over $M$ and define $E_1, \dots, E_n$ as the orthonormal frame dual to $\omega_i$:
\begin{equation}
    E_i\omega_j = \delta_{ij}
\end{equation}
Then, there exist $a_1, \dots, a_n$ such that
\begin{equation}
    X = \sum_i a_i E_i
\end{equation}
It's easy to see that
\begin{equation}
    g(X, E_i) = g(\sum a_i E_i, E_i) = \sum a_i g(E_i, E_i) = a_i
\end{equation}
And since $g$ and $E_i$ are smooth, that means $a_i$ are smooth as well. So, define the bijection $f: TM \rightarrow T^*M$ as the function the for each $(p, X)$ gives
\begin{equation}
    f(p, X) = \sum_i a_i \omega_i
\end{equation}
Since $a_i$ are smooth and $\omega_i$ are smooth, this is a smooth bijection with a smooth inverse (defined the same way), and therefore a diffeomorphism.

\section*{Question 4}
Let $M \subset \R^n$ be a $n-1$-dimensional smooth manifold and let $\omega = dx_1 \wedge \cdots \wedge dx_n$ be the standard volume form on $\R^n$. Let $N$ be the unit normal vector field. Then $\iota_N \omega$ is a $n-1$-form over $M$. Define an orthonormal frame $e_1, \dots, e_{n-1}$ over $M$. Since $\langle N, X \rangle = 0$ for every $X \in TM$, the set $N, e_1, \dots, e_{n-1}$ is orthonormal and thus a basis for $\R^n$.
Then:
\begin{equation}
    \iota_N \omega(e_1, \dots, e_{n-1}) = \omega(N, e_1, \dots, e_{n-1}) = \pm 1
\end{equation}
Which means $\iota_N \omega$ is a Riemannian volume form on $M$.

\section*{Question 5}
Let $X, Y, X \in \Gamma(TM)$ be vector fields and let $\nabla$ be a torsion free metric compatible connection. Then, because it is metric compatible:
\begin{equation}
\begin{gathered}
    X g(Y,Z) = g(\nabla_X Y, Z) + g(Y, \nabla_X Z)
    \\
    Y g(Z,X) = g(\nabla_Y Z, X) + g(Z, \nabla_Y X)
    \\
    Z g(X,Y) = g(\nabla_Z X, Y) + g(X, \nabla_Z Y)
\end{gathered}
\end{equation}
Adding the first two equations and substracting the last one:
\begin{equation}
\begin{gathered}
    X g(Y,Z) + Y g(Z,X) + Z g(X,Y) = \\ =
    g(\nabla_X Y, Z) + g(Y, \nabla_X Z) + g(\nabla_Y Z, X) + g(Z, \nabla_Y X) - g(\nabla_Z X, Y) - g(X, \nabla_Z Y)
\end{gathered}
\end{equation}
But since $\nabla$ is torsion free:
\begin{equation}
\begin{gathered}
    g(Z, \nabla_Y X) = g(Z, \nabla_X Y) - g(Z, [X, Y])
    \\
    g(Y, \nabla_X Z) - g(\nabla_Z X, Y) = g(Y, [X, Z])
    \\
    g(\nabla_Y Z, X) - g(X, \nabla_Z Y) = g(X, [Y, Z])
\end{gathered}
\end{equation}
This gives:
\begin{equation}
    X g(Y,Z) + Y g(Z,X) - Z g(X,Y) = 2g(Z, \nabla_X Y) - g(Z, [X, Y]) + g(Y, [X, Z]) + g(X, [Y, Z])
\end{equation}
Reogranizing:
\begin{equation}
    2g(Z, \nabla_X Y) = X g(Y,Z) + Y g(Z,X) - Z g(X,Y) + g(Z, [X, Y]) - g(Y, [X, Z]) - g(X, [Y, Z])
\end{equation}\footnote{The worksheet has a wrong sign here, but I double checked and this is the formulation. I actually didn't notice it until I did Q11, where terms vanished with this formula wrong}

\section*{Question 6}
Let $V$ be a vector field along a curve $\gamma(t)$ and let $\tilde{V}_1$ and $\tilde{V}_2$ be two extensions of the vector field to a small neighborhood around the curve. Then, for every connection $\nabla$, using the linearity of the connection:
\begin{equation}
    \nabla_{\dot{\gamma}(t)} \left(\tilde{V}_1 - \tilde{V}_2\right) = 
    \nabla_{\dot{\gamma}(t)} \tilde{V}_1 - \nabla_{\dot{\gamma}(t)} \tilde{V}_2
\end{equation}
But since along the curve $\tilde{V}_1 = \tilde{V}_2 = V$, then $\tilde{V}_1 - \tilde{V}_2 = 0$ which gives
\begin{equation}
    \nabla_{\dot{\gamma}(t)} \tilde{V}_1 - \nabla_{\dot{\gamma}(t)} \tilde{V}_2 = 0
\end{equation}
\begin{equation}
    \nabla_{\dot{\gamma}(t)} \tilde{V}_1 = \nabla_{\dot{\gamma}(t)} \tilde{V}_2
\end{equation}
But since this is just the covariant derivative of $V$, it means the covariant derivation is independent of the extension.

\section*{Question 7}
Let $\nabla$ be the Levi Civita connection, i.e. the torsion free metric compatible connection. Then, define a local coordinate system $x_1, \dots, x_n$ for every point $p \in M$ and a local frame $\partial_{x_1},\dots, \partial_{x_n}$. Then, the Christoffel symbols are
\begin{equation}
    \nabla_{\partial_i} \partial_j = \sum_n \Gamma^{n}_{ij} \partial_n
\end{equation}
Applying the metric with $\partial_l$:
\begin{equation}
    g\left(\partial_l, \nabla_{\partial_i} \partial_j\right) = g\left(\partial_l, \sum_n \Gamma^{n}_{ij} \partial_n\right)
\end{equation}
The RHS gives:
\begin{equation}
    g\left(\partial_l, \sum_n \Gamma^{n}_{ij} \partial_n\right) = \sum_n \Gamma^{n}_{ij} g(\partial_l, \partial_n) = g_{ln}\Gamma^{n}_{ij}
\end{equation}
Where the last step is written in Einstein summation convention. Multiplying by the inverse metric $g^{lk}$ and summing over $l$ gives:
\begin{equation}
    g^{lk} g_{ln}\Gamma^{n}_{ij} = \delta_{kn}\Gamma^{n}_{ij} = \Gamma^{k}_{ij}
\end{equation}
And so:
\begin{equation}
    \Gamma^{k}_{ij} = g^{lk} g\left(\partial_l, \nabla_{\partial_i} \partial_j\right)
\end{equation}
Using Koszul formula that we proved in Question 5:
\begin{equation}
    \Gamma^{k}_{ij} = \frac{1}{2}g^{lk} \left[
        \partial_i g(\partial_j,\partial_l) +
        \partial_j g(\partial_l,\partial_i) -
        \partial_l g(\partial_i,\partial_j) +
        g(\partial_l, [\partial_i, \partial_j]) -
        g(\partial_j, [\partial_i, \partial_l]) -
        g(\partial_i, [\partial_j, \partial_l])
    \right]
\end{equation}
But since $[\partial_i, \partial_j] = 0$ for every $i, j$:
\begin{equation}
    \Gamma^{k}_{ij} = \frac{1}{2}g^{lk} \left[
        \partial_i g_{jl} +
        \partial_j g_{il} -
        \partial_l g_{ij}
    \right]
\end{equation}

\section*{Question 8}
Let $M = S^2$ be isometrically embedded in $\R^3$ and let $g$ be the Eucledian metric. Now, let $\gamma(t)$ be the equator. In $\R^3$ Eucledian coordinates:
\begin{equation}
    \gamma(t) = (\cos 2\pi t, \sin 2\pi t, 0)
\end{equation}
Now, let $X = \partial_x$ be a parallel vector field in $\R^3$ and let $Y$ be its orthogonal projection. Then on $\gamma(t)$ it can calculted by projecting on the tangent to the unit circle and then back to the Eucledian coordinates:
\begin{equation}
    Y|_{\gamma(t)} = -\sin 2\pi t (\sin 2\pi t, \cos 2\pi t, 0)
\end{equation}
So, for $t = 0$, $t = 0.125$ and $t = 0.25$ the projected vector field is:
\begin{equation}
    Y|_{\gamma(0)} = (0, 0, 0)
    \qquad
    Y|_{\gamma(0.125)} = (0.5, 0.5, 0)
    \qquad
    Y|_{\gamma(0.25)} = (1, 0, 0)
\end{equation}
So, this is an example of an orthogonal projection of a parall vector field that is not parallel in an isometrically embedded manifold in $\R^n$.

\section*{Question 11}
Let $(M, g)$ be an oriented Riemannian surface and $(e_1, e_2)$ a local oriented orthonormal frame. Let $X \in TM$ be a vector field and let $\nabla$ be the Levi Civita connection. writing $X$ in terms of $(e_1, e_2)$:
\begin{equation}
    X = x_1 e_1 + x_2 e_2
\end{equation}
Gives:
\begin{equation}
    \nabla_X e_1 = x_1 \nabla_{e_1} e_1 + x_2 \nabla_{e_2} e_1
    \qquad
    \nabla_X e_2 = x_1 \nabla_{e_1} e_2 + x_2 \nabla_{e_2} e_2
\end{equation}
The object $\nabla_{e_i} e_j$ can be calcualted by applying the metric with $e_k$ to it, and using Koszul formula:
\begin{equation}
    g(\nabla_{e_i} e_j, e_k) = \frac{1}{2}\left[
        e_i g(e_j,e_k) + e_j g(e_k,e_i) - e_k g(e_i,e_j) + g(e_k, [e_i, e_j]) - g(e_j, [e_i, e_k]) - g(e_i, [e_j, e_k])
    \right]
\end{equation}
And using $e_i g(e_j, e_k) = e_i g_{jk} = e_i \delta_{jk} = 0$:
\begin{equation}
    g(\nabla_{e_i} e_j, e_k) = \frac{1}{2}\left[g(e_k, [e_i, e_j]) - g(e_j, [e_i, e_k]) - g(e_i, [e_j, e_k])
    \right]
\end{equation}
Writing:
\begin{equation}
    \nabla_{e_i} e_j = a_{ij} e_1 + b_{ij} e_2
\end{equation}
Gives:
\begin{equation}
    a_{ij} = g(\nabla_{e_i} e_j, e_1)
    \qquad
    b_{ij} = g(\nabla_{e_i} e_j, e_2)
\end{equation}
Specifically:
\begin{equation}
\begin{gathered}
    a_{11} = g(\nabla_{e_1} e_1, e_1) = \frac{1}{2}\left[g(e_1, [e_1, e_1]) - g(e_1, [e_1, e_1]) - g(e_1, [e_1, e_1])
    \right] = 0
    \\
    a_{12} = g(\nabla_{e_1} e_2, e_1) = \frac{1}{2}\left[g(e_1, [e_1, e_2]) - g(e_2, [e_1, e_1]) - g(e_1, [e_2, e_1])
    \right] = g(e_1, [e_1, e_2])
    \\
    a_{21} = g(\nabla_{e_2} e_1, e_1) = \frac{1}{2}\left[g(e_1, [e_2, e_1]) - g(e_1, [e_2, e_1]) - g(e_2, [e_1, e_1])
    \right] = 0
    \\
    a_{22} = g(\nabla_{e_2} e_2, e_1) = \frac{1}{2}\left[g(e_1, [e_2, e_2]) - g(e_2, [e_2, e_1]) - g(e_2, [e_2, e_1])
    \right] = g(e_2, [e_1, e_2])
\end{gathered}
\end{equation}
And similarly:
\begin{equation}
    b_{12} = b_{22} = 0
    \qquad
    b_{11} = g(e_1, [e_2, e_1])
    \qquad
    b_{21} = g(e_2, [e_2, e_1])
\end{equation}
So, plugging those into the above formulation:
\begin{equation}
\begin{gathered}
    \nabla_{e_1} e_1 = g(e_1, [e_2, e_1]) e_2
    \qquad
    \nabla_{e_1} e_2 = g(e_1, [e_1, e_2]) e_1
    \qquad
    \nabla_{e_2} e_1 = g(e_2, [e_2, e_1]) e_2
    \qquad
    \nabla_{e_2} e_2 = g(e_2, [e_1, e_2]) e_1
\end{gathered}
\end{equation}
And:
\begin{equation}
\begin{gathered}
    \nabla_X e_1 = x_1 g(e_1, [e_2, e_1]) e_2 + x_2 g(e_2, [e_2, e_1]) e_2
    \\
    \nabla_X e_2 = x_1 g(e_1, [e_1, e_2]) e_1 + x_2 g(e_2, [e_1, e_2]) e_1
\end{gathered}
\end{equation}
So, defining the form
\begin{equation}
    \omega = g(e_1, [e_2, e_1]) \omega_1 + g(e_2, [e_2, e_1]) \omega_2
\end{equation}
where $\omega_1, \omega_2$ be the coframe dual to $e_1, e_2$ gives:
\begin{equation}
    \nabla e_1 = \omega \otimes e_2
    \qquad
    \nabla e_2 = - \omega \otimes e_1
\end{equation}
This is a very similar structure to the cannoical connection $\varphi$. Then there exist a 1-form $\varphi$ such that: 
\begin{equation}
    d\omega_1 = \varphi \wedge \omega_2
    \qquad
    d\omega_2 = -\varphi \wedge \omega_1
\end{equation}
And then
\begin{equation}
    d\varphi = K \omega_1 \wedge \omega_2
\end{equation}
But we notice that:
\begin{equation}
    \omega \wedge \omega_2 = g(e_1, [e_2, e_1]) \omega_1 \wedge \omega_2
\end{equation}
But also, since $e_j(\omega_1 e_i) = e_j\delta_{i1} = 0$:
\begin{equation}
    d\omega_1(e_i, e_j) = e_i \omega_1 e_j + e_j \omega_1 e_i - \omega_1 ([e_i, e_j]) = \omega_1 ([e_j, e_i])
\end{equation}
Which gives $d\omega_1(e_1, e_1) = d\omega_1(e_2, e_2) = 0$ and $d\omega_1(e_1, e_2) = -d\omega_1(e_2, e_1) = \omega_1 ([e_2, e_1]) = g(e_1, [e_2, e_1])$. This gives:
\begin{equation}
    d\omega_1 = g(e_1, [e_2, e_1]) \omega_1 \wedge \omega_2 = \omega \wedge \omega_2
\end{equation}
A similar calculation for $\omega \wedge \omega_1$ and $d\omega_2$ will give the opposite equation. So, $\omega$ follows the first structure equation, and thus $\omega = \varphi$ and
\begin{equation}
    d\omega = K \omega_1 \wedge \omega_2
\end{equation}
\end{document}